\begin{pastexam}[label={px:ncu-114-dsa-08}]{中央資工所 114 學年度 第 8 題}
（複選；每題 5 分，全部答對才給分，答錯倒扣 1 分）

（第 8--10 題共用下面的說明和程式。）

Consider a UAV (Unmanned Aerial Vehicle) with F unit of fuel travels from NCU to a destination
``target'' km away. There are n gas stations along the way. When the UAV refuels at a gas station, all
the fuel of the gas station is transferred into the UAV.

Suppose that the UAV consumes one unit of fuel for every kilometer it travels. Let the position and fuel
of a gas station indicate the distance between the gas station and NCU and its volume of fuel,
respectively. Return the minimum number of refueling stops the UAV must make in order to reach its
destination.
\begin{minted}[fontsize=\small]{c}
typedef struct {
    int position;
    int fuel;
} stop_info;

int Refuel(int target, int F, int n, stop_info s[]) {
    int N_R = 0;   // number of refuel
    int i;   //gas stations ID
    stop_info X;

    while (F < target) {
        for (i=0;  L1  ; ++i)
            Q.push(s[i].fuel);
        if (Q.empty()) return -1;
        X=Q.pop();
        F += X.fuel
        N_R++;
    }
    return N_R;
}
\end{minted}

What should be filled in blank \textbf{L1}?
\begin{enumerate}[label=(\alph*)]
\item (i<n) \&\& (s[i].position <=F)
\item (i<n) \&\& (s[i].fuel <=F)
\item (i<n) \&\& (s[i].position <=target)
\item i<n
\end{enumerate}
\end{pastexam}

\begin{pastsoln}
答案：(a)。

\textbf{先把程式讀對}。 原卷的程式有幾處筆誤：\texttt{F += X.fuel} 少了分號；\texttt{Q} 推入的是
\texttt{s[i].fuel}（整數），彈出的卻當成 \texttt{stop\_info} 用；\texttt{i} 在每一輪 \texttt{while}
都從 $0$ 重新開始，會把同一個加油站重複推入、之後可能被重複使用。 題意是下面這個貪心法（第 8--10 題
都照這個讀法）：加油站依 position 由小到大排好；\texttt{i} 只在進迴圈前設成 $0$，之後一路往後走，每個
加油站只推入一次；\texttt{Q} 存油量，\texttt{pop} 傳回目前最大的油量。
\begin{pseudo}{最少加油次數（題目的貪心法）}
\Function{Refuel}{$\mathit{target}, F, s[1..n]$}
    \State $r \gets 0$；$i \gets 1$；$Q \gets$ 空的最大堆積
    \While{$F < \mathit{target}$}
        \While{$i \le n$ \AND $s[i].\mathit{position} \le F$} \Comment{目前油量到得了的站}
            \State \Call{Insert}{$Q, s[i].\mathit{fuel}$}；$i \gets i + 1$
        \EndWhile
        \If{$Q$ 是空的}
            \State \Return $-1$ \Comment{到不了}
        \EndIf
        \State $F \gets F + \Call{Extract-Max}{Q}$；$r \gets r + 1$
    \EndWhile
    \State \Return $r$
\EndFunction
\end{pseudo}
$F$ 在這裡的意思是「出發以來總共拿到的油」，也就是目前最遠飛得到的位置（每公里耗一單位油）。 經過
一個站不一定要停；等到油不夠時，再「事後」決定在已經經過的站裡挑油最多的那一站加油。 這樣做是對的：
已經飛過的站，當時停下來加油和事後補算，能到達的最遠位置一樣，所以每次不夠時挑油最多的站，加油次數
最少（資結書第~19.1.6~節的貪心選擇性質：任何最佳解都可以換成先用油最多的那一站，次數不增加）。

\textbf{L1 的作用}：把「目前到得了」而且還沒推入的站推進 \texttt{Q}。 到得了的條件是
$\mathit{position} \le F$。
\begin{itemize}
\item (a) \texttt{(i<n) \&\& (s[i].position <=F)}：正是上面的條件，而且因為站依位置排好，遇到第一個
到不了的站就停，後面的站下一輪再看。 對。
\item (b) 比的是油量和 $F$，和到不到得了無關。 例如 $F = 10$、唯一的站在位置 $5$、油量 $20$：到得了，
卻因為 $20 > 10$ 不推入，回傳 $-1$。 錯。
\item (c) 只要在終點之前的站都推入，包括目前飛不到的站。 例如 $\mathit{target} = 100$、$F = 10$、
唯一的站在位置 $50$、油量 $100$：飛不到那一站，正確答案是 $-1$，(c) 卻會在那裡加油而回傳 $1$。 錯。
\item (d) 同 (c)，所有的站都推入。 錯。
\end{itemize}
\end{pastsoln}
